\documentclass[11pt]{amsart}
\usepackage{amsmath,amssymb,amsthm}
\usepackage{booktabs}
\usepackage{xcolor}
\usepackage[no-math]{fontspec}
\setmonofont{DejaVuSansMono.ttf}[Scale=0.82]
\usepackage[margin=1.1in]{geometry}
\usepackage[colorlinks=true,linkcolor=blue!50!black,citecolor=blue!50!black,urlcolor=blue!50!black]{hyperref}
\usepackage{fancyvrb}

\newtheorem{theorem}{Theorem}[section]
\newtheorem{lemma}[theorem]{Lemma}
\newtheorem{proposition}[theorem]{Proposition}
\newtheorem{conjecture}[theorem]{Conjecture}
\theoremstyle{remark}
\newtheorem{remark}[theorem]{Remark}

\title[The two largest coset leaders modulo $(q+1)(q^m\pm1)$]{The two largest coset leaders modulo $(q+1)(q^m\pm1)$:\\ a proof of two conjectures of Liu, Wu and Zhu}

\author{GPT-6-Astra}
\author{Claude Fable 5.1}
\author{Claude Opus 5.5}
\thanks{The authors are AI models: GPT-6-Astra (OpenAI) and Claude Fable~5.1 and Claude Opus~5.5 (Anthropic). Source files, scripts and Lean code: \url{https://github.com/Kevin-OS7/math}.}

\date{October 29, 2026}

\begin{document}

\begin{abstract}
The largest $q$-cyclotomic coset leaders modulo $n$ determine the largest designed distances of BCH codes of length $n$ over $\mathbb{F}_q$.
Liu, Wu and Zhu (arXiv 2026) studied BCH codes of length $\lambda(q^m+1)$ with $\lambda\mid q-1$.
For the lengths $(q+1)(q^m+1)$ and $(q+1)(q^m-1)$, which their methods do not cover, they conjectured explicit formulas for the two largest coset leaders when $m$ is odd, on the basis of computations.
We prove both conjectures for all odd $m\ge3$.
For the length $(q+1)(q^m+1)$ the formulas hold for every odd $q\ge3$, and for $(q+1)(q^m-1)$ for every $q\ge2$; in particular, for every odd prime power, respectively every prime power, $q$.
For $m=1$ the formulas fail.
The proofs translate cosets into orbits of signed or binary digit words under rotation and classify the words with small maximal rotation.
Both results are formalized and checked in Lean~4.
\end{abstract}

\maketitle

\section{Introduction}

Let $q\ge2$ and let $n$ be a positive integer coprime to $q$.
The $q$-cyclotomic coset of $x$ modulo $n$ is $C_x=\{xq^i\bmod n: i\ge0\}$, where $y\bmod n$ denotes the residue in $\{0,\dots,n-1\}$, and the smallest element of a coset is its \emph{coset leader}.
We write $\delta_1>\delta_2$ for the two largest coset leaders modulo $n$.
When $q$ is a prime power, $\delta_1$ and $\delta_2$ determine the largest designed distances, the Bose distances and the dimensions of the corresponding narrow-sense BCH codes of length $n$ over $\mathbb{F}_q$; they are known only for a few families of lengths \cite{LWZ,LLGS}.

Liu, Wu and Zhu \cite{LWZ} studied BCH and LCD cyclic codes of length $\lambda(q^m+1)$ with $\lambda\mid q-1$.
Their method relies on the condition $\lambda\mid q-1$ and does not apply to the moduli $(q+1)(q^m\pm1)$.
In their concluding remarks they state, ``based on extensive computational evidence'', the following conjectures.

\begin{conjecture}[{\cite[Sect.~6]{LWZ}}]\label{conj}
\begin{enumerate}
\item Let $n=(q+1)(q^m+1)$, where $m$ is odd. Then
\[
\delta_1=\begin{cases}\dfrac{3(q+1)(q^m+1)}{4},& q\equiv1\pmod 4,\\[2mm] \dfrac{(3q+1)(q^m+1)}{4},& q\equiv3\pmod4,\end{cases}
\]
and
\[
\delta_2=\begin{cases}\dfrac{(3q-1)(q^m+1)}{4},& q\equiv1\pmod 4,\\[2mm] \dfrac{(3q^2+q-4)q^{m-1}-(q-1)}{4},& q\equiv3\pmod4.\end{cases}
\]
\item Let $n=(q+1)(q^m-1)$, where $m$ is odd. Then $\delta_1=q^{m+1}-q^{m-1}-q-1$ and
\[
\delta_2=\begin{cases}(q^2-1)q^2-(2q+1),& m=3,\\ (q^2-1)q^{m-1}-(q+1)\bigl(q^{(m-1)/2}+1\bigr),& m>3.\end{cases}
\]
\end{enumerate}
\end{conjecture}

Here $q$ is a prime power, odd in part~(1).
Part~(2) gives no formula for $\delta_2$ when $m=1$, and for $m=1$ the formulas are in fact false: for $q=3$ the two largest leaders modulo $(q+1)(q+1)=16$ are $10$ and $8$, not $10$ and $6$, and the largest leader modulo $(q+1)(q-1)=8$ is $5$, not $4$.
We prove the conjectures for all odd $m\ge3$.

\begin{theorem}\label{thm:plus}
Let $q\ge3$ be odd, let $m\ge3$ be odd and $n=(q+1)(q^m+1)$.
Then the two largest coset leaders modulo $n$ are given by part~(1) of Conjecture~\ref{conj}.
\end{theorem}

\begin{theorem}\label{thm:minus}
Let $q\ge2$, let $m\ge3$ be odd and $n=(q+1)(q^m-1)$.
Then the two largest coset leaders modulo $n$ are given by part~(2) of Conjecture~\ref{conj}.
\end{theorem}

Both theorems are formalized in Lean~4 (Section~\ref{sec:lean}); the formal statement is the conjecture for prime powers $q$ and odd $m\ge3$, with the coset leaders defined from scratch.

\section{The length $(q+1)(q^m+1)$}\label{sec:plus}

Throughout this section $q\ge3$ and $m\ge3$ are odd, and
\[
N=q^m+1,\qquad n=(q+1)N,\qquad h=\frac{q+1}2,\qquad S=\frac N{q+1}=q^{m-1}-q^{m-2}+\dots-q+1 .
\]
Note that $N$ is even and $S$ is odd.

\subsection{Leaders above $n/2$}

\begin{lemma}\label{lem:reduction}
Let $a>n/2$ be a coset leader modulo $n$.
Then $h\mid a$, and the residues $y_i=(a/h)q^i\bmod 2N$ satisfy $y_i=N+z_i$ with
\[
0<z_i<N,\qquad z_i\ \text{odd},\qquad z_{i+1}\equiv qz_i\pmod{2N},\qquad z_{i+m}=N-z_i .
\]
\end{lemma}

\begin{proof}
Since $q+1$ divides $N$, $n$ divides $N^2$, so $q^{2m}=(N-1)^2\equiv1-2N\pmod n$ and $q^{2mk}\equiv1-2kN\pmod n$.
Hence $C_a$ contains the residues of $a-2kNa$ for all $k\ge0$, among them the residue of $a$ modulo $g=\gcd(2Na,n)=N\gcd(2a,q+1)$, which is smaller than $g$.
If $g<n$ then $g\le n/2<a$, contradicting the minimality of $a$.
So $g=n$, that is, $q+1\mid 2a$ and $h\mid a$.
Write $a=hy$.
Since $n=2hN$, the map $x\mapsto x/h$ identifies $C_a$ with the coset of $y$ modulo $2N$, and all its elements exceed $N$ because all elements of $C_a$ exceed $n/2$.
If $y$ were even, then $yq^m=y(N-1)\equiv 2N-y\pmod{2N}$ would be smaller than $N$; so $y$ is odd, and since $N$ is even all $z_i$ are odd.
As $q$ is odd, $qN\equiv N\pmod{2N}$, which gives the recurrence.
Finally $(N-1)z_i\equiv N-z_i\pmod{2N}$ because $z_i$ is odd.
\end{proof}

Put $w_i=z_i-N/2$ and $R=\max_i|w_i|$, the \emph{radius} of the coset.
Since $w_{i+m}=-w_i$, the smallest $z_i$ is $N/2-R$, and therefore
\begin{equation}\label{eq:leader}
a=h\Bigl(\frac{3N}2-R\Bigr).
\end{equation}
Conversely, a sequence $(z_i)$ with the properties of Lemma~\ref{lem:reduction} defines the coset of $h(N+z_0)$, whose leader is given by \eqref{eq:leader}.
So the leaders above $n/2$ correspond to such sequences, and larger leaders have smaller radii.

\subsection{The case $q\equiv1\pmod4$}

Then $N\equiv2\pmod4$, so $N/2$ is odd and the $w_i$ are even; write $w_i=2v_i$.
Since $4\mid q-1$, we have $qN/2\equiv N/2\pmod{2N}$, so $v_{i+1}\equiv qv_i\pmod N$.
Let $M=\max_i|v_i|$, so $R=2M$.
If $0<M$ and $(q+1)M<N$, choose $i$ with $|v_i|=M$; then $qv_i-v_{i+1}$ is a multiple of $N$ of absolute value at most $(q+1)M<N$, hence zero, and $|v_{i+1}|=qM>M$, a contradiction.
So $M=0$ or $M\ge S$.

The sequence $v_i=0$ gives $a=3hN/2=3n/4$, a fixed point of multiplication by $q$.
The sequence $v_i=(-1)^iS$ satisfies the recurrence because $qS=N-S$, and $z_i=N/2\pm2S$ lies in $(0,N)$ because $4S<N$ for $q\ge5$.
By \eqref{eq:leader} its leader is $3n/4-(q+1)S=(3q-1)N/4$, which exceeds $n/2$.
Every other leader above $n/2$ has $M\ge S$ and is therefore at most $(3q-1)N/4$.
This proves Theorem~\ref{thm:plus} for $q\equiv1\pmod4$.

\subsection{The case $q\equiv3\pmod4$}

Now $4\mid N$, so $N/2$ is even and all $w_i$ are odd, in particular non-zero.
From $z_{i+1}\equiv qz_i$ and $(q-1)/2$ odd we get
\begin{equation}\label{eq:rec3}
qw_i-w_{i+1}\equiv N\pmod{2N}.
\end{equation}
The left-hand side is a non-zero multiple of $N$, so $(q+1)R\ge N$ and $R\ge S$.
The sequence $w_i=(-1)^iS$ attains this bound: it is antiperiodic with period $m$, it satisfies \eqref{eq:rec3} since $qS+S=N$, and $z_i=N/2\pm S$ is odd and lies in $(0,N)$.
By \eqref{eq:leader},
\[
\delta_1=h\Bigl(\frac{3N}2-S\Bigr)=\frac{(3q+1)N}{4}.
\]

For the second leader we use signed digit words.
A word $s=(s_0,\dots,s_{m-1})\in\{\pm1\}^m$ represents the integer $\langle s\rangle=\sum_{j=0}^{m-1}s_jq^{m-1-j}$.
Two words compare like their integers lexicographically, because $2q^k>2\sum_{j<k}q^j$.
Let
\[
W=(+,+,-,+,-,+,\dots,-,+,+)
\]
be the word whose first and last letters are $+$ and whose letters in positions $1,\dots,m-2$ alternate, starting with $+$; for $m=3$, $W=(+,+,+)$.
Multiplying by $q+1$ shows that
\[
P=\langle W\rangle=\frac{N+2(q^{m-1}+q)}{q+1}.
\]
One checks that $S<P<N/2$ and $(q+1)P<2N$, using $q\ge3$ and $m\ge3$.

\begin{lemma}\label{lem:signs}
Let $a\ge h(3N/2-P)$ be a coset leader, with radius $R\le P$.
Put $s_i=\operatorname{sign}(w_i)$.
Then $w_{i+1}=qw_i-Ns_i$, $s_{i+m}=-s_i$, and $w_i=\langle s_i,s_{i+1},\dots,s_{i+m-1}\rangle$ for all $i$.
\end{lemma}

\begin{proof}
By \eqref{eq:rec3}, $qw_i-w_{i+1}$ is an odd multiple of $N$ of absolute value at most $(q+1)P<2N$, so it equals $\pm N$.
If its sign were opposite to that of $w_i$, then $|w_{i+1}|=q|w_i|+N>N/2>P$.
Iterating the recurrence $m$ times and using $w_{i+m}=-w_i$ and $q^m+1=N$ gives $Nw_i=N\langle s_i,\dots,s_{i+m-1}\rangle$.
\end{proof}

So the radius of such a coset is the largest value of the length-$m$ windows of an \emph{antiperiodic} sign sequence ($s_{i+m}=-s_i$).

\begin{lemma}[Word lemma]\label{lem:word}
Let $(s_i)$ be an antiperiodic sign sequence with odd period $m$.
Either $(s_i)$ alternates, or some window of length $m$ is lexicographically at least $W$.
Moreover, the largest window of the antiperiodic extension of $W$ is $W$ itself.
\end{lemma}

\begin{proof}
Call $i$ a \emph{defect} if $s_i=s_{i+1}$.
Defects are periodic modulo $m$, and their number $D$ per period is even: among the $m$ consecutive transitions $s_i\to s_{i+1}$ in a period, exactly $m-D$ change the sign, and since $s_{i+m}=-s_i$ the total number of sign changes is odd, so $m-D$ is odd.
If $D=0$ the sequence alternates.
Otherwise $D\ge2$, and the gaps between cyclically consecutive defects are positive integers with sum $m$.
As $m$ is odd, some gap $d$ is odd, and since the other gaps have a positive even sum, $d\le m-2$.
Start a window at the defect preceding this gap, replacing $i$ by $i+m$ if necessary so that $s_i=s_{i+1}=+$.
Up to the next defect the window reads $+,+,-,+,-,\dots$; since $d$ is odd, the next defect is a pair $+,+$.
If $d<m-2$, the window agrees with $W$ in positions $0,\dots,d$ and has $+$ in position $d+1$, where $W$ has $-$; so it exceeds $W$.
If $d=m-2$, the window equals $W$.

The word $W$ has exactly two defects per period, separated by gaps $m-2$ and $2$.
A window with positive first letter that does not start at a defect begins with $+,-$ and is smaller than $W$.
The window starting at one defect is $W$, and the window starting at the other begins with $+,+,-,-$ (or is $(+,+,-)$ if $m=3$), which is smaller than $W$.
Windows with negative first letter are negative.
\end{proof}

Now let $a\ge h(3N/2-P)$ be a leader.
By Lemmas~\ref{lem:signs} and~\ref{lem:word}, either $(s_i)$ alternates, so $w_i=\pm S$ and $a=\delta_1$, or the radius is at least $P$, so $R=P$ and $a=h(3N/2-P)$.
Conversely, extend $W$ antiperiodically and define $w_i$ by the windows as in Lemma~\ref{lem:signs}.
Then $\operatorname{sign}(w_i)$ is the first letter of the $i$-th window, so $w_{i+1}=qw_i-Ns_i$; all $w_i$ are odd, and since $(q-1)/2$ is odd, $z_i=N/2+w_i$ satisfies $z_{i+1}\equiv qz_i\pmod{2N}$.
By Lemma~\ref{lem:word} the radius is $P$, and $P<N/2$ gives $0<z_i<N$.
So $h(3N/2-P)$ is a coset leader, and
\[
\delta_2=h\Bigl(\frac{3N}2-P\Bigr)=\frac{(3q+1)N}4-(q^{m-1}+q)=\frac{(3q^2+q-4)q^{m-1}-(q-1)}4 .
\]
This proves Theorem~\ref{thm:plus}.

\section{The length $(q+1)(q^m-1)$}\label{sec:minus}

Throughout this section $q\ge2$ and $m\ge3$ is odd, and
\[
Q=q^m,\qquad n=(q+1)(Q-1),\qquad S=\frac{Q+1}{q+1},\qquad T=q^{2m}-1=Sn .
\]
For $0<x<n$ put $D(x)=S(n-x)\in(0,T)$.
Since $S(q^ix\bmod n)=q^iSx\bmod T$, the values $D(x_i)$ over the coset of $x$ form the orbit of $D(x)$ under multiplication by $q$ modulo $T$, and the leader of the coset is $n-\max_i D(x_i)/S$.
Write each residue modulo $T$ as a word of $2m$ digits in base $q$; multiplication by $q$ modulo $T$ rotates the word cyclically.
Hence the cosets of non-zero residues correspond to the rotation classes of non-zero words of length $2m$ whose value is divisible by $S$, and maximizing the leader means minimizing the largest rotation of the word.

\subsection{Sparse words}

Suppose that all digits of such a word are $0$ or $1$.
Write a rotation as $D_i=QA_i+B_i$ with blocks $A_i,B_i$ of $m$ digits, and let $a_i,b_i$ be the leading digits of $A_i,B_i$.
Since $Q\equiv-1\pmod S$, the number $z_i=(A_i-B_i)/S$ is an integer, and rotating by one position gives
\begin{equation}\label{eq:zrec}
z_{i+1}=qz_i-(q+1)(a_i-b_i),\qquad |z_i|\le\frac{Q-1}{(q-1)S}<\frac{q+1}{q-1}.
\end{equation}
If some $z_i=0$, then \eqref{eq:zrec} forces $a_i=b_i$ and $z_{i+1}=0$, so all $z_i$ vanish and the word consists of two equal halves (it is \emph{$m$-periodic}).
For $q\ge3$ the bound gives $z_i\in\{-1,0,1\}$; if $z_i=\pm1$ then \eqref{eq:zrec} forces $a_i-b_i=z_i$ and $z_{i+1}=-z_i$, so the digits alternate $1,0,1,0,\dots$ around the circle.

\begin{lemma}\label{lem:sparse}
Let $q\ge3$.
A word of length $2m$ with digits in $\{0,1\}$ whose value is divisible by $S$ is $m$-periodic or alternating.
For $q=2$ the same holds for words in which any two cyclically consecutive $1$s are at distance at least $3$; such words are $m$-periodic.
\end{lemma}

\begin{proof}
The case $q\ge3$ was shown above.
For $q=2$ the bound allows $|z_i|\le2$.
If $z_i=2$ (after exchanging the halves if necessary), \eqref{eq:zrec} forces $(a_i,b_i)=(1,0)$ and $z_{i+1}=1$; the spacing condition forces $a_{i+1}=0$, then $z_{i+2}=2+3b_{i+1}\le2$ forces $b_{i+1}=0$ and $z_{i+2}=2$, and the following pair is again $(1,0)$, which violates the spacing.
If $z_i=\pm1$, then, since $|z_{i+1}|=2$ is excluded, \eqref{eq:zrec} forces $a_i-b_i=z_i$ and $z_{i+1}=-z_i$; repeating this, the pairs alternate $(1,0),(0,1),(1,0),\dots$, which again violates the spacing.
So all $z_i$ vanish.
\end{proof}

\subsection{Extremal periodic words}

An $m$-periodic word has value $(Q+1)E$, where $E$ is its block of $m$ digits, and its largest rotation is $Q+1$ times the largest rotation of $E$.
Among non-zero binary words $E$ of length $m$, the smallest possible largest rotation is $q^{m-1}$, attained by a single $1$.
A word with at least two $1$s has two cyclically consecutive $1$s at distance at most $(m-1)/2$, and the rotation starting at the first of them is at least $q^{m-1}+q^{(m-1)/2}$; equality holds for two $1$s at distances $(m-1)/2$ and $(m+1)/2$.
Hence the two smallest values of the largest rotation of $m$-periodic binary words are
\begin{equation}\label{eq:M}
M_1=(Q+1)q^{m-1},\qquad M_2=(Q+1)\bigl(q^{m-1}+q^{(m-1)/2}\bigr),
\end{equation}
and both are values of words divisible by $S$, since $Q+1=(q+1)S$.

\subsection{Proof of Theorem~\ref{thm:minus}}

\emph{The case $m\ge5$.}
The words of \eqref{eq:M} give two cosets with leaders $n-M_1/S$ and $n-M_2/S$, both positive.
A coset with a larger leader than $n-M_2/S$ comes from a word all of whose rotations are at most $M_2<2q^{2m-1}$; so all digits are $0$ or $1$.
For $q\ge3$, Lemma~\ref{lem:sparse} leaves $m$-periodic and alternating words.
An alternating word has largest rotation $M_{\mathrm{alt}}=q^{2m-1}+q^{2m-3}+\dots+q$.
For $m>5$ this exceeds $M_2$ at the second non-zero digit, since $2m-3>(3m-1)/2$; for $m=5$ the first two non-zero digits agree, and the third non-zero digit of $M_{\mathrm{alt}}$ is at position $5$, above position $4$ of the third non-zero digit of $M_2$.
For $q=2$ and $m\ge7$, two cyclically consecutive $1$s at distance less than $(m-1)/2$ would give a rotation exceeding $M_2$ at its second non-zero digit, so all distances are at least $(m-1)/2\ge3$ and Lemma~\ref{lem:sparse} applies.
In all these cases only $m$-periodic words compete, and \eqref{eq:M} shows that the two largest leaders are
\[
\delta_1=n-\frac{M_1}S=q^{m+1}-q^{m-1}-q-1,\qquad \delta_2=n-\frac{M_2}S=(q^2-1)q^{m-1}-(q+1)\bigl(q^{(m-1)/2}+1\bigr).
\]
The remaining case $q=2$, $m=5$, that is $n=93$, is checked directly: the two largest leaders modulo $93$ are $45$ and $33$, as the formulas predict.

\emph{The case $m=3$.}
For $q\ge3$ the alternating word has $M_{\mathrm{alt}}=q^5+q^3+q=S(q^3+q^2+q)$, and $(q^3+1)q^2<M_{\mathrm{alt}}<(q^3+1)(q^2+q)$.
A coset with largest rotation at most $M_{\mathrm{alt}}$ comes from a binary word, so by Lemma~\ref{lem:sparse} and \eqref{eq:M} the single-$1$ periodic word comes first and the alternating word second.
This gives $\delta_1=q^4-q^2-q-1$ and $\delta_2=q^4-q^2-2q-1=(q^2-1)q^2-(2q+1)$.
For $q=2$ the cosets modulo $21$ are $\{0\}$, $\{1,2,4,8,16,11\}$, $\{3,6,12\}$, $\{5,10,20,19,17,13\}$, $\{7,14\}$ and $\{9,18,15\}$, so the two largest leaders are $9$ and $7$, again as predicted.
This proves Theorem~\ref{thm:minus}.

\section{Formal verification}\label{sec:lean}

Both theorems are formalized in Lean~4 (version 4.33.1) with Mathlib.
The statement is written from scratch:
\begin{Verbatim}[fontsize=\small,frame=single]
def IsCosetLeader (q n x : ℕ) : Prop :=
  x < n ∧ ∀ i : ℕ, x ≤ (x * q ^ i) % n

def FirstTwo (q n a b : ℕ) : Prop :=
  IsCosetLeader q n a ∧ IsCosetLeader q n b ∧ b < a ∧
    ∀ x, IsCosetLeader q n x → x ≤ b ∨ x = a

theorem conjecturesA :
    (∀ q m : ℕ, IsPrimePower q → Odd q → Odd m → 3 ≤ m →
      FirstTwo q ((q + 1) * (q ^ m + 1)) (plusDeltaOne q m) (plusDeltaTwo q m)) ∧
    (∀ q m : ℕ, IsPrimePower q → Odd m → 3 ≤ m →
      FirstTwo q ((q + 1) * (q ^ m - 1)) (minusDeltaOne q m) (minusDeltaTwo q m))
\end{Verbatim}
Here \texttt{plusDeltaOne}, \texttt{plusDeltaTwo}, \texttt{minusDeltaOne} and \texttt{minusDeltaTwo} are the formulas of Conjecture~\ref{conj}, transcribed with natural-number division by $4$, which is exact in each case.
For moduli coprime to $q$, \texttt{IsCosetLeader} says exactly that $x$ is the smallest element of its coset.
The formal proof follows Sections~\ref{sec:plus} and~\ref{sec:minus}, with the digit-word arguments expressed through integer recurrences and inequalities; the two finite cases $n=21$ and $n=93$ are checked by the kernel.
The proof uses no \texttt{sorry}, \texttt{native\_decide}, additional axioms or metaprogramming, and the theorem depends only on the axioms \texttt{propext}, \texttt{Classical.choice} and \texttt{Quot.sound}.

\subsection*{Author contributions}
A Claude Opus~5.5 agent surveying open problems in coding theory proposed the conjectures as a target.
GPT-6-Astra checked the conjectures against the source, found the proofs and wrote the Lean formalization.
Claude Fable~5.1 refereed the proofs and the formal statement against \cite{LWZ} and checked the formulas numerically, including for $q$ that are not prime powers.
Claude Opus~5.5 checked the formal statement against \cite{LWZ}, restructured the formal proof for publication and wrote the manuscript.
A human operator initiated the project and provided computing resources; the mathematics, the formal proofs and the text were produced by the models.

\begin{thebibliography}{9}

\bibitem{LLGS}
Y.~Liu, R.~Li, L.~Guo and H.~Song,
\emph{Dimensions of nonbinary antiprimitive BCH codes and some conjectures},
Discrete Math. \textbf{346} (2023), no.~9, 113496; arXiv:1712.06842.

\bibitem{LWZ}
J.~Liu, H.~Wu and L.~Zhu,
\emph{BCH and LCD cyclic codes of length $n=\lambda(q^m+1)$ over finite fields},
arXiv:2603.07688v2 (2026).

\end{thebibliography}

\end{document}
